\begin{proof}[Proof of \cref{obj:lem:tensor-jackson-coefficient-envelope}]
Write
\[
\mathbb T_{\mathcal J}:=[-\pi,\pi]^{\mathcal J},
\qquad
N_K:=2(K-1),
\]
and define the tensor convolution
\[
(\mathcal J_K f)(\varphi)
:=
\int_{\mathbb T_{\mathcal J}}
f\{\cos(\varphi-\tau)\}
\prod_{\jmath\in\mathcal J}J_K(\tau_\jmath)\,d\tau .
\]
For the positive order \(K\) in the lemma, write \(J_K=c_KG_K\), where
\[
G_K(t)=
\begin{cases}
\left\{\dfrac{\sin(Kt/2)}{\sin(t/2)}\right\}^{4},& \sin(t/2)\ne0,\\[1ex]
K^4,& \sin(t/2)=0,
\end{cases}
\]
with the removable values filled continuously, and set
\[
c_K:=\left(\int_{-\pi}^{\pi}G_K(t)\,dt\right)^{-1}.
\]
Then \(G_K(t)\ge0\) for every \(t\), and \(G_K(0)=K^4>0\), so \(\int_{-\pi}^{\pi}G_K(t)\,dt>0\) and \(c_K>0\).  Thus \(J_K\ge0\) and
\[
\int_{-\pi}^{\pi}J_K(t)\,dt
=
c_K\int_{-\pi}^{\pi}G_K(t)\,dt
=1 .
\]
Fubini's theorem for the nonnegative continuous tensor product then gives
\[
\prod_{\jmath\in\mathcal J}J_K(\tau_\jmath)\ge 0,
\qquad
\int_{\mathbb T_{\mathcal J}}
\prod_{\jmath\in\mathcal J}J_K(\tau_\jmath)\,d\tau=1 .
\]
Continuity of \(f\) on the normalized cube and continuity of the coordinatewise cosine map give the integrability needed for the displayed convolutions and for the linearity identities below. % lean: tensorJackson_nonneg, tensorJackson_integral_eq_one

\begin{enumerate}
\item
First construct the polynomial representation.  The Jackson frequency bound follows from the Chebyshev form of the raw kernel.  If \(U_{K-1}\) denotes the Chebyshev polynomial of the second kind, then
\[
G_K(s)=\{U_{K-1}(\cos(s/2))\}^{4}.
\]
Since \(U_{K-1}(-\varrho)=(-1)^{K-1}U_{K-1}(\varrho)\), the fourth power is an even polynomial in \(\varrho\).  Hence there is a polynomial \(Q_K^{\mathrm{even}}\) of degree at most \(2(K-1)=N_K\) such that
\[
\{U_{K-1}(\varrho)\}^{4}=Q_K^{\mathrm{even}}(\varrho^2).
\]
Using \(\cos^2(s/2)=(1+\cos s)/2\), this gives
\[
G_K(s)=Q_K^{\mathrm{even}}\{(1+\cos s)/2\}.
\]
Multiplication by the constant \(c_K\) preserves the frequency bound, and a polynomial of degree at most \(N_K\) in \(\cos s\) is a trigonometric polynomial with frequencies at most \(N_K\).  Thus there are real coefficient sequences \(\chi^{\cos}_{K,r}\) and \(\chi^{\sin}_{K,r}\), \(0\le r\le N_K\), such that
\[
J_K(s)=\sum_{r=0}^{N_K}
\{\chi^{\cos}_{K,r}\cos(rs)+\chi^{\sin}_{K,r}\sin(rs)\}.
\]
% lean: jackson_isTrigPolyLE
Fix all coordinates except \(\omega_0\in\mathcal J\).  For an angle vector \(\eta:\mathcal J\to\mathbb R\) and \(t\in\mathbb R\), write
\[
(\eta_{\omega_0}\leftarrow t)_\omega
=
\begin{cases}
t,&\omega=\omega_0,\\
\eta_\omega,&\omega\ne\omega_0.
\end{cases}
\]
Set, for \(\psi\in\mathbb T_{\mathcal J}\),
\[
W_{\omega_0,\eta}(\psi)
:=
f(\cos \psi)
\prod_{\omega\in\mathcal J\setminus\{\omega_0\}}
J_K(\eta_\omega-\psi_\omega).
\]
The period-box translation \(\psi=(\eta_{\omega_0}\leftarrow t)-\tau\), with the box understood modulo the common \(2\pi\)-period in each coordinate, gives
\[
(\mathcal J_K f)(\eta_{\omega_0}\leftarrow t)
=
\int_{\mathbb T_{\mathcal J}}
f(\cos\psi)
\prod_{\omega\in\mathcal J}
J_K\{(\eta_{\omega_0}\leftarrow t)_\omega-\psi_\omega\}\,d\psi
=
\int_{\mathbb T_{\mathcal J}}
W_{\omega_0,\eta}(\psi)J_K(t-\psi_{\omega_0})\,d\psi .
\]
The translation is valid because the displayed integrand is continuous and is \(2\pi\)-periodic in each coordinate; the periodicity follows from the cosine term and the trigonometric-polynomial representation of \(J_K\). % lean: tensorConvolution_shifted
Define
\[
\Xi^{\cos}_{\omega_0,\eta,r}
:=
\int_{\mathbb T_{\mathcal J}}
W_{\omega_0,\eta}(\psi)
\{\chi^{\cos}_{K,r}\cos(r\psi_{\omega_0})-\chi^{\sin}_{K,r}\sin(r\psi_{\omega_0})\}\,d\psi,
\]
\[
\Xi^{\sin}_{\omega_0,\eta,r}
:=
\int_{\mathbb T_{\mathcal J}}
W_{\omega_0,\eta}(\psi)
\{\chi^{\cos}_{K,r}\sin(r\psi_{\omega_0})+\chi^{\sin}_{K,r}\cos(r\psi_{\omega_0})\}\,d\psi .
\]
Substituting the finite expansion of \(J_K(t-\psi_{\omega_0})\), using the angle-subtraction identities for sine and cosine, and integrating the finite sum term by term gives
\[
(\mathcal J_K f)(\eta_{\omega_0}\leftarrow t)
=
\sum_{r=0}^{N_K}
\{\Xi^{\cos}_{\omega_0,\eta,r}\cos(rt)+\Xi^{\sin}_{\omega_0,\eta,r}\sin(rt)\}.
\]
Thus each one-coordinate section is a trigonometric polynomial of frequency at most \(N_K\).  The same section is even in that coordinate: reflecting the integration variable \(\tau_{\omega_0}\mapsto-\tau_{\omega_0}\) preserves \(\mathbb T_{\mathcal J}\), the kernel is even, and cosine is even.  Therefore the sine part cancels.  Interpolating the resulting even trigonometric polynomial in one cosine coordinate at the nodes
\[
\xi_j=\frac{j}{N_K+1},\qquad
\zeta_j=\arccos(\xi_j),\qquad j=0,\ldots,N_K,
\]
and repeating this coordinate by coordinate gives a polynomial \(\mathcal R\) in the four variables such that
\[
\mathcal R(\cos\varphi)=(\mathcal J_K f)(\varphi)
\quad\text{for every }\varphi:\mathcal J\to\mathbb R.
\]
The interpolation basis has degree at most \(N_K\) in the coordinate being added, while the induction hypothesis controls the other coordinates.  Hence, for every
\(\beta\in\operatorname{supp}(\mathcal R)\),
\[
\beta_\omega\le N_K=2(K-1)
\quad\text{for every }\omega\in\mathcal J,
\]
and
\[
\deg(\mathcal R)\le 4N_K=8(K-1).
\]
% lean: tensorConvolution_exists_mvPolynomial

\item
The convolution is bounded by the same envelope \(B\).  Since
\(\lvert f(z)\rvert\le B\) on the normalized cube,
\[
B-f(z)\ge0,
\qquad
B+f(z)\ge0
\quad
(z\in[-1,1]^{\mathcal J}).
\]
Kernel nonnegativity gives
\[
(\mathcal J_K(B-f))(\varphi)\ge0,
\qquad
(\mathcal J_K(B+f))(\varphi)\ge0 .
\]
By linearity of the integral and the unit-mass identity,
\[
(\mathcal J_K(B-f))(\varphi)
=
B-(\mathcal J_K f)(\varphi),
\qquad
(\mathcal J_K(B+f))(\varphi)
=
B+(\mathcal J_K f)(\varphi).
\]
Therefore
\[
-B\le(\mathcal J_K f)(\varphi)\le B,
\qquad
\lvert(\mathcal J_K f)(\varphi)\rvert\le B
\quad
(\varphi:\mathcal J\to\mathbb R).
\]
% lean: hsub, hadd, hconv_bound

\item
This angle-space bound transfers to the whole normalized cube.  For any
\(z\in[-1,1]^{\mathcal J}\), take
\[
\varphi_z(\omega):=\arccos(z_\omega),
\qquad \omega\in\mathcal J,
\]
so that \(\cos(\varphi_z)_\omega=z_\omega\).  The representation from the first step and the bound from the second step give
\[
|\mathcal R(z)|
=
|\mathcal R(\cos\varphi_z)|
=
|(\mathcal J_K f)(\varphi_z)|
\le B .
\]
% lean: hp_bound

\item
It remains to bound the coefficient one-norm.  The coefficient estimate used here is the following.  Let \(\mathcal S\) be a polynomial in \(\delta_{\mathrm{dim}}\) normalized variables, each coordinate degree at most \(N_K\), and suppose
\[
|\mathcal S(y)|\le B_{\mathrm{env}}
\qquad
(y\in[-1,1]^{\delta_{\mathrm{dim}}}),
\qquad
B_{\mathrm{env}}\ge0 .
\]
Then
\[
\sum_{\rho\in\operatorname{supp}(\mathcal S)}
|\mathcal S_\rho|
\le
\{2(N_K+1)^2 3^{N_K}\}^{\delta_{\mathrm{dim}}}
B_{\mathrm{env}} .
\]
For one variable, if \(s\) has degree at most \(N_K\) and \(|s(x)|\le C\) on \([-1,1]\), then \(s(\cos t)\) is an even trigonometric polynomial,
\[
s(\cos t)=\sum_{r=0}^{N_K} a_r\cos(rt).
\]
Orthogonality on \([-\pi,\pi]\) gives
\[
|a_r|\le 2C
\qquad (0\le r\le N_K).
\]
Since \(\cos(rt)=T_r(\cos t)\) and the Chebyshev coefficient one-norm satisfies
\[
\|T_r\|_1\le 3^r,
\]
the polynomial identity \(s(x)=\sum_{r=0}^{N_K}a_rT_r(x)\) yields
\[
\|s\|_1
\le
\sum_{r=0}^{N_K}|a_r|\,\|T_r\|_1
\le
2(N_K+1)3^{N_K}C .
\]
For the multivariate estimate, write
\[
\mathcal S(x_0,y)
=
\sum_{r=0}^{N_K}s_r(y)x_0^r .
\]
For each fixed \(y\in[-1,1]^{\delta_{\mathrm{dim}}}\), the preceding one-dimensional bound gives
\[
|s_r(y)|
\le
2(N_K+1)3^{N_K}B_{\mathrm{env}} .
\]
Applying the induction hypothesis to each coefficient polynomial \(s_r\) and summing over \(r=0,\ldots,N_K\) gives
\[
\|\mathcal S\|_1
=
\sum_{r=0}^{N_K}\|s_r\|_1
\le
(N_K+1)\{2(N_K+1)^2 3^{N_K}\}^{\delta_{\mathrm{dim}}}
\{2(N_K+1)3^{N_K}B_{\mathrm{env}}\},
\]
which is exactly
\[
\|\mathcal S\|_1
\le
\{2(N_K+1)^2 3^{N_K}\}^{\delta_{\mathrm{dim}}+1}
B_{\mathrm{env}} .
\]
The case \(\delta_{\mathrm{dim}}=0\) is the constant-polynomial case and follows directly from the displayed uniform bound. % lean: mvCoeffL1_le_bounded

\item
Apply this coefficient estimate to \(\mathcal S=\mathcal R\), \(\delta_{\mathrm{dim}}=4\), and \(B_{\mathrm{env}}=B\).  With \(N_K=2(K-1)\),
\[
\sum_{\beta\in\operatorname{supp}(\mathcal R)}
|\mathcal R_\beta|
\le
\{2(N_K+1)^2 3^{N_K}\}^{4}B .
\]
The elementary bound
\[
N_K+1\le 2^{N_K+1}
\]
and \(3^{N_K}\le4^{N_K}\) imply
\[
2(N_K+1)^2 3^{N_K}
\le
2\{2^{N_K+1}\}^2 4^{N_K}
=
2^{4N_K+3}.
\]
Consequently
\[
\{2(N_K+1)^2 3^{N_K}\}^{4}
\le
2^{4(4N_K+3)} .
\]
Since \(N_K=2(K-1)\) and \(K\ge1\),
\[
4(4N_K+3)
=
32(K-1)+12
\le
40K+20 .
\]
Because the base \(2\) is at least \(1\) and \(B\ge0\),
\[
\sum_{\beta\in\operatorname{supp}(\mathcal R)}
|\mathcal R_\beta|
\le
2^{40K+20}B .
\]
Together with the representation, coordinate-support bound, and total-degree bound proved above, this is the asserted polynomial \(\mathcal R\). % lean: tensorConvolution_exists_mvPolynomial_four_coeffBound
\end{enumerate}
\end{proof}
